\documentclass[UTF8,a4paper,11pt]{ctexart}

\usepackage{geometry}
\usepackage{hyperref}
\usepackage{longtable}
\usepackage{array}
\usepackage{enumitem}
\usepackage{xcolor}
\usepackage{fancyvrb}

\geometry{left=2.4cm,right=2.4cm,top=2.4cm,bottom=2.4cm}
\hypersetup{
  colorlinks=true,
  linkcolor=blue,
  urlcolor=blue,
  citecolor=blue
}
\setlist[itemize]{leftmargin=2em}
\setlist[enumerate]{leftmargin=2em}
\DefineVerbatimEnvironment{code}{Verbatim}{fontsize=\small}

\title{\textbf{EyeProtect 项目技术总述}\\\large 算法依据、工具调用与实现边界}
\author{EyeProtect Project}
\date{2026 年 7 月}

\begin{document}
\maketitle
\tableofcontents
\newpage

\section{项目定位}

EyeProtect 是一个部署在 Android 平板端的护眼原型项目。当前版本已经完成：

\begin{itemize}
  \item 通过 Android 系统接口读取屏幕/App 使用时间；
  \item 通过 CameraX 调用前置摄像头；
  \item 通过 Google ML Kit Face Detection 检测人脸和左右眼位置；
  \item 绘制左右眼关键点、眼部边界框，并输出眼部坐标；
  \item 计算 FPS、检测耗时、人脸框、眼部中心点；
  \item 裁剪并显示左右眼 ROI；
  \item 在后台、锁屏、横竖屏场景下管理摄像头资源。
\end{itemize}

当前项目的眼部识别目标是“眼部位置检测”，不是虹膜检测、眼球方向判断或屏幕注视点预测。

\section{总体技术架构}

项目由 Android 原生能力、CameraX 视频采集、ML Kit 人脸检测和自定义规则算法组成。

\begin{code}
Android UsageStatsManager
  -> 屏幕/App 使用时间

CameraX Preview
  -> 前置摄像头实时预览

CameraX ImageAnalysis
  -> 实时读取视频帧
  -> InputImage
  -> Google ML Kit Face Detection
  -> Face / FaceContour / FaceLandmark
  -> 左右眼点位
  -> 边界框、中心点、平滑、ROI 裁剪
  -> Overlay 绘制与 Logcat 输出
\end{code}

\section{功能与依据}

\subsection{屏幕使用时间统计}

屏幕使用时间统计基于 Android 平台 API：\texttt{UsageStatsManager}。

实现依据：

\begin{itemize}
  \item Android 系统会记录应用进入前台的使用情况；
  \item 应用获得“使用情况访问权限”后，可以查询一段时间内的应用使用统计；
  \item 项目使用 \texttt{queryUsageStats} 查询今日统计，并读取 \texttt{totalTimeInForeground} 作为使用时长。
\end{itemize}

当前代码位置：

\begin{itemize}
  \item \texttt{app/src/main/java/com/example/eyeprotect/usage/UsageAccessHelper.kt}
  \item \texttt{app/src/main/java/com/example/eyeprotect/usage/ScreenUsageRepository.kt}
\end{itemize}

算法流程：

\begin{code}
检查 Usage Access 权限
  -> 未授权：跳转系统设置页
  -> 已授权：UsageStatsManager.queryUsageStats(...)
  -> 按 packageName 聚合 totalTimeInForeground
  -> 按使用时长倒序展示
\end{code}

这部分不是视觉算法，也不是后台录屏；它是系统级使用统计读取。

\subsection{前置摄像头调用}

前置摄像头调用基于 AndroidX CameraX。CameraX 是 Android Jetpack 中的相机库，用于简化不同 Android 设备上的相机兼容性处理。

项目中使用两个 CameraX use case：

\begin{itemize}
  \item \texttt{Preview}：负责把摄像头画面显示到界面；
  \item \texttt{ImageAnalysis}：负责实时读取帧，送入眼部检测链路。
\end{itemize}

关键实现：

\begin{code}
CameraSelector.DEFAULT_FRONT_CAMERA
ImageAnalysis.STRATEGY_KEEP_ONLY_LATEST
ImageAnalysis.Builder().setTargetResolution(Size(640, 480))
\end{code}

当前代码位置：

\begin{itemize}
  \item \texttt{app/src/main/java/com/example/eyeprotect/camera/FrontCameraController.kt}
\end{itemize}

选择 \texttt{STRATEGY\_KEEP\_ONLY\_LATEST} 的原因是实时检测更重视低延迟，不需要处理积压旧帧。

\subsection{人脸与眼部识别}

眼部识别当前基于 Google ML Kit Face Detection SDK。

项目中使用的依赖：

\begin{code}
implementation("com.google.mlkit:face-detection:16.1.7")
\end{code}

需要特别说明：当前 Android 版本不是使用 Python/OpenCV/MediaPipe 实现的。Python 版本曾作为 PC 端验证方案存在，但平板端实际运行的是 CameraX + ML Kit。

ML Kit Face Detection 在项目中的作用：

\begin{itemize}
  \item 检测单帧图像中的人脸；
  \item 输出人脸框；
  \item 输出左右眼 landmark；
  \item 输出左右眼 contour，即眼部轮廓点。
\end{itemize}

关键配置：

\begin{code}
FaceDetectorOptions.Builder()
  .setPerformanceMode(FaceDetectorOptions.PERFORMANCE_MODE_FAST)
  .setLandmarkMode(FaceDetectorOptions.LANDMARK_MODE_ALL)
  .setContourMode(FaceDetectorOptions.CONTOUR_MODE_ALL)
  .setClassificationMode(FaceDetectorOptions.CLASSIFICATION_MODE_ALL)
  .enableTracking()
  .build()
\end{code}

当前代码位置：

\begin{itemize}
  \item \texttt{app/src/main/java/com/example/eyeprotect/vision/EyePositionAnalyzer.kt}
\end{itemize}

\subsection{眼部点位提取}

ML Kit 返回 \texttt{Face} 对象后，项目优先读取：

\begin{code}
FaceContour.LEFT_EYE
FaceContour.RIGHT_EYE
\end{code}

如果 contour 不可用，则使用：

\begin{code}
FaceLandmark.LEFT_EYE
FaceLandmark.RIGHT_EYE
\end{code}

并基于中心点生成一个近似眼框。

算法逻辑：

\begin{code}
Face
  -> getContour(LEFT_EYE).points
  -> getContour(RIGHT_EYE).points
  -> 如果轮廓点为空，回退到 getLandmark(...)
\end{code}

\subsection{眼部边界框计算}

眼部边界框不是 ML Kit 直接给出的眼框，而是项目根据眼部点位计算出来的。

设眼部点集为：

\begin{code}
P = [(x1, y1), (x2, y2), ..., (xn, yn)]
\end{code}

边界框计算：

\begin{code}
box.x1 = min(x) - padding
box.y1 = min(y) - padding
box.x2 = max(x) + padding
box.y2 = max(y) + padding
\end{code}

之后裁剪到图像范围：

\begin{code}
x1 >= 0
y1 >= 0
x2 <= imageWidth - 1
y2 <= imageHeight - 1
\end{code}

当前代码位置：

\begin{itemize}
  \item \texttt{EyePositionAnalyzer.calculateBoundingBox(...)}
  \item \texttt{app/src/main/java/com/example/eyeprotect/vision/BoundingBox.kt}
\end{itemize}

\subsection{眼部中心点}

眼部中心点基于边界框计算：

\begin{code}
centerX = (box.x1 + box.x2) / 2
centerY = (box.y1 + box.y2) / 2
\end{code}

注意：当前中心点表示“眼部框中心”，不是虹膜中心，也不是视线落点。

\subsection{坐标平滑}

实时检测中的眼部点会轻微抖动。项目使用指数平滑降低视觉抖动。

公式：

\begin{code}
smoothed = previous + alpha * (current - previous)
\end{code}

当前 \texttt{alpha} 设置为 0.35。该值越大，响应越快但抖动越明显；该值越小，画面更稳但跟随延迟更高。

当前代码位置：

\begin{itemize}
  \item \texttt{EyePositionAnalyzer.smoothEye(...)}
\end{itemize}

\subsection{左右眼 ROI 裁剪}

ROI 是 Region of Interest，即感兴趣区域。项目根据平滑后的眼部边界框，从当前帧图像中裁剪左右眼小图。

流程：

\begin{code}
CameraX ImageProxy
  -> YUV_420_888
  -> NV21
  -> JPEG
  -> Bitmap
  -> 按旋转角度修正
  -> 根据 eyeBox 扩大 1.8 倍
  -> Bitmap.createBitmap(...)
\end{code}

当前代码位置：

\begin{itemize}
  \item \texttt{app/src/main/java/com/example/eyeprotect/vision/ImageProxyBitmapConverter.kt}
  \item \texttt{EyePositionAnalyzer.cropEyeRoi(...)}
\end{itemize}

ROI 当前只用于显示和后续算法准备；尚未用于虹膜识别或眼球方向判断。

\section{使用的工具、SDK 与开源边界}

\begin{longtable}{|p{3.2cm}|p{3.4cm}|p{5.1cm}|p{3.2cm}|}
\hline
\textbf{组件} & \textbf{用途} & \textbf{项目中的调用方式} & \textbf{开源/开放边界} \\
\hline
Android Platform API & 系统使用统计、权限、生命周期 & \texttt{UsageStatsManager}、\texttt{AppOpsManager}、\texttt{BroadcastReceiver} & Android 平台 API，AOSP 中有大量平台代码；设备实现可能由厂商定制。 \\
\hline
AndroidX CameraX & 前置摄像头预览与帧分析 & \texttt{Preview}、\texttt{ImageAnalysis}、\texttt{CameraSelector.DEFAULT\_FRONT\_CAMERA} & AndroidX/Jetpack 开源库。 \\
\hline
Google ML Kit Face Detection & 人脸、眼部 landmark/contour 检测 & \texttt{FaceDetection.getClient(...)}、\texttt{FaceContour.LEFT\_EYE}、\texttt{FaceContour.RIGHT\_EYE} & 公开 SDK；底层模型不是本项目可修改的开源模型。 \\
\hline
Kotlin + Gradle & Android 应用开发与构建 & Kotlin 源码、Gradle Android Plugin、依赖管理 & Kotlin 与 Gradle 为开源工具链。 \\
\hline
自定义规则算法 & 边界框、中心点、平滑、ROI & \texttt{calculateBoundingBox}、\texttt{smoothEye}、\texttt{cropEyeRoi} & 项目自研规则代码，可完全修改。 \\
\hline
\end{longtable}

\section{ML Kit 与 MediaPipe 的关系说明}

项目早期曾实现过一个 PC 端 Python 版本：

\begin{code}
Python + OpenCV + MediaPipe Face Mesh
\end{code}

该版本用于验证眼部点位检测思路。但当前可安装到华为平板的 Android APK 没有运行 Python、OpenCV 或 MediaPipe Python 代码。

当前 Android APK 的实际链路是：

\begin{code}
CameraX
  -> ML Kit Face Detection
  -> FaceContour.LEFT_EYE / RIGHT_EYE
  -> 自定义边界框与 ROI 算法
\end{code}

因此，正式描述应写为：

\begin{quote}
本项目的 Android 端眼部位置识别基于 Google ML Kit Face Detection SDK。系统通过 CameraX 获取前置摄像头实时视频帧，将帧输入 ML Kit Face Detector，利用其输出的人脸框、左右眼轮廓点和眼部关键点，计算左右眼边界框、中心点和 ROI 区域。
\end{quote}

\section{当前输出数据}

项目当前输出的核心数据包括：

\begin{itemize}
  \item \texttt{timestamp}：时间戳；
  \item \texttt{face\_id}：人脸跟踪 ID；
  \item \texttt{face\_count}：检测到的人脸数量；
  \item \texttt{fps}：平滑后的帧率；
  \item \texttt{latency\_ms}：单帧检测耗时；
  \item \texttt{source\_width/source\_height}：分析帧尺寸；
  \item \texttt{face\_box}：人脸框；
  \item \texttt{left\_eye\_box/right\_eye\_box}：左右眼框；
  \item \texttt{left\_eye\_center/right\_eye\_center}：左右眼中心点。
\end{itemize}

示例：

\begin{code}
{
  "timestamp": 1783651200125,
  "face_id": 0,
  "face_count": 1,
  "fps": 24.8,
  "latency_ms": 37,
  "source_width": 640,
  "source_height": 480,
  "face_box": {"x1": 180, "y1": 70, "x2": 460, "y2": 390},
  "left_eye_box": {"x1": 245, "y1": 160, "x2": 305, "y2": 195},
  "left_eye_center": {"x": 275, "y": 178},
  "right_eye_box": {"x1": 340, "y1": 160, "x2": 400, "y2": 195},
  "right_eye_center": {"x": 370, "y": 178}
}
\end{code}

\section{资源与生命周期策略}

摄像头是高占用资源，项目在工程层面实现了以下策略：

\begin{itemize}
  \item 应用进入后台时释放摄像头；
  \item 应用回前台时，如果此前摄像头正在运行，则自动恢复；
  \item 屏幕关闭或锁屏时停止摄像头，不自动恢复；
  \item 用户可以点击 \texttt{Stop Camera} 手动释放摄像头；
  \item 横竖屏切换时保存恢复意图，避免无谓丢失状态。
\end{itemize}

当前代码位置：

\begin{itemize}
  \item \texttt{MainActivity.onPause}
  \item \texttt{MainActivity.onResume}
  \item \texttt{MainActivity.ACTION\_SCREEN\_OFF BroadcastReceiver}
  \item \texttt{FrontCameraController.stopPreview}
\end{itemize}

\section{当前限制}

当前版本仍有明确边界：

\begin{itemize}
  \item 不能检测虹膜中心；
  \item 不能判断眼球向左/向右/向上/向下；
  \item 不能预测屏幕注视点；
  \item 眼部中心点是眼框中心，不是瞳孔或虹膜中心；
  \item ML Kit Face Detection 的底层模型不可在本项目内训练或修改；
  \item 在弱光、遮挡、强侧脸、多人脸场景下检测稳定性会下降。
\end{itemize}

\section{后续扩展建议}

建议下一步顺序：

\begin{enumerate}
  \item 睁闭眼检测：使用 ML Kit 的 \texttt{leftEyeOpenProbability/rightEyeOpenProbability}；
  \item 眨眼事件检测：实现 \texttt{OPEN -> CLOSED -> OPEN} 状态机；
  \item 疲劳提醒：结合闭眼时长、眨眼频率、屏幕使用时间；
  \item 距离过近提醒：基于人脸框高度比例或双眼中心距离；
  \item 虹膜与视线方向：评估 MediaPipe Face Landmarker / Iris 能力或独立模型；
  \item 屏幕注视点预测：需要用户校准数据，不建议在无校准条件下直接实现。
\end{enumerate}

\section{参考资料}

\begin{itemize}
  \item Android \texttt{UsageStatsManager} 官方文档：\url{https://developer.android.com/reference/android/app/usage/UsageStatsManager}
  \item Android CameraX 图像分析官方文档：\url{https://developer.android.com/media/camera/camerax/analyze}
  \item AndroidX Camera 发布说明：\url{https://developer.android.com/jetpack/androidx/releases/camera}
  \item Google ML Kit Face Detection Android 文档：\url{https://developers.google.com/ml-kit/vision/face-detection/android}
  \item Google ML Kit Face Detection Concepts：\url{https://developers.google.com/ml-kit/vision/face-detection}
  \item Kotlin 官方网站：\url{https://kotlinlang.org/}
  \item Gradle 官方网站：\url{https://gradle.org/}
\end{itemize}

\end{document}
